\section{pacing 提案への反応が継続}
Anthropic のスローダウン提案に対し、Hassabis は方向性を支持し業界横断の標準団体案と接続した。Musk は競争他社によるピアレビュー型監視を支持し、Artificial Analysis も第三者アクセス拡大を歓迎した。
\dailyxsource{Demis Hassabis (@demishassabis)}{https://x.com/demishassabis/status/2098909516582490602}

\section{Jacob Coxon 辞職・警告の二次波}
元 OpenAI / Anthropic 研究者 Jacob Coxon の辞職と強い安全警告が、ABC や WIRED などの窓内投稿で再拡散した。辞職と本人スレッド自体は窓外で、窓内ではメディア二次波として観測された。
\dailyxsource{ABC (@ABC)}{https://x.com/ABC/status/2099222596273230069}

\section{Ethan Mollick: Astra と Fable 5.1 の経済影響}
Mollick は、Astra と Fable 5.1 が適切に誘導されれば数週間分の人的作業を安定して担えるとし、pacing があっても経済変化は避けられないと主張した。能力評価は本人の観測で未検証。
\dailyxsource{Ethan Mollick (@emollick)}{https://x.com/emollick/status/2099235643792462040}

\section{Yann LeCun がオープンソース危険論を批判}
LeCun は、Dario が2019年時点で GPT-2 のオープンソース化を危険視していたことを引き合いに出し、当時も今も批判的だと投稿した。「escapes」言説も重大な過失か規制キャプチャ目的だと主張した。
\dailyxsource{Yann LeCun (@ylecun)}{https://x.com/ylecun/status/2099248236074545576}

\section{オープンソース規制・AIカルテル懸念}
スローダウン提案を、オープンソース禁止や計算資源不足の言い換えと見る投稿が広く共有された。DeepSeek 4.1 の性能主張と結び付け、競争抑制だと解釈する声もあったが、政策意図を裏付ける材料は不足している。
\dailyxsource{Tazerface16 (@Tazerface16)}{https://x.com/Tazerface16/status/2099089300164088274}

\section{NCP / Next Concept Prediction 論文が拡散}
解説アカウントが、8.9B規模の NCP-ArchPreview と「次トークンではなく離散概念を予測する」方式を紹介し拡散した。数学性能向上や計算削減などの数値は解説投稿経由で、Grok収集では一次検証していない。
\dailyxsource{HowToPrompt (@HowToPrompt\_\_)}{https://x.com/HowToPrompt__/status/2099186312033259653}

\section{Recurrent Looped Transformer への批判}
無限推論深度をうたう Recurrent Looped Transformer の投稿がバイラル化したことを受け、Microsoft の研究者が「実験が一つもない論文が拡散している」と批判した。論文内容は本収集では精読していない。
\dailyxsource{Alexia Jolicoeur-Martineau (@jm\_alexia)}{https://x.com/jm_alexia/status/2099136049158988101}

\section{BRICS オープンソースAI表明と偽情報}
習近平が BRICS 向けにオープンソースAIや技術訓練を約束したとする投稿が流通した。一方、「中国が米国研究の減速を支持」とする画像付き投稿も拡散し、別投稿者が実際はBRICS OSS推進だと訂正した。
\dailyxsource{Eng China (@Eng\_china5)}{https://x.com/Eng_china5/status/2099048399282811031}

\section{DeepSeek-V4.1-Flash のローカル推論報告}
ローカル推論ラボが、DeepSeek V4.1 Flash を4基の RTX 6000 Pro 上で高速推論したと自己報告した。別投稿では拒否回路を除いた改変ウェイトも主張された。速度・安全性数値はいずれも未検証。
\dailyxsource{Your Local AI Lab (@YourLocalAILab)}{https://x.com/YourLocalAILab/status/2099201400555073890}

\section{Musk: Nvidia AI computers を宇宙へ}
Musk は、翌年 SpaceX が「Nvidia VR NLV72 AI computers」を宇宙へ打ち上げることに高い自信があると投稿した。製品名や契約に対応する公式一次資料は本収集では確認されていない。
\dailyxsource{Elon Musk (@elonmusk)}{https://x.com/elonmusk/status/2099255050581295238}

\section{Meta Muse のエージェント先行フィード}
Microsoft 所属を名乗る投稿者が、Muse の個人情報に基づくエージェント先行フィードを高く評価し、プロンプト不要で先回りするUIが大規模普及に必要だと述べた。機能の公式一次情報は未確認。
\dailyxsource{nicbstme (@nicbstme)}{https://x.com/nicbstme/status/2099211882628596216}

\section{OpenAI IPO 報道が再循環}
Reuters が、Altman が安全リスクを理由に2026年内のIPOを否定したとする内容を窓内で再掲した。コミュニティでは日程や意図をめぐる別解釈も出たが、それらは投稿者側の見方である。
\dailyxsource{Reuters (@Reuters)}{https://x.com/Reuters/status/2099241738518892731}

\section{Hugging Face 侵入事案の事後解釈}
窓外の Hugging Face 侵入事案を、pacing 提案の動機として読み直す投稿が続いた。過失か意図か、報酬ハッキング、キルスイッチ法案の実効性など複数の解釈があるが、窓内は事後議論が中心。
\dailyxsource{Yann LeCun (@ylecun)}{https://x.com/ylecun/status/2099087198918197654}

\section{3B未満の小型オープンウェイトLLM需要}
OpenAI 所属を名乗る投稿者が、ブラウザや拡張機能で動く3B未満の最良オープンウェイトLLMを問い、反応を集めた。特定モデルの優劣ではなく、オンデバイス推論への実務需要を示す観測として記録する。
\dailyxsource{VB (@reach\_vb)}{https://x.com/reach_vb/status/2099140289717833842}

